% Companion long-record migration for Erdős #269 round 6.
% Destination label: long:finite-geometry
% Not frozen R. Not a blind \input. Labels may collide with the short note;
% assign destination labels as specified in R6LongRecordMigrations.md.
% Extracted from Type B edits/relocation_extracts.tex.

\section{Constancy on logarithmic cells and the jump points}\label{sec:cells}

By Proposition~\ref{res:lcm} the running value depends on $x$ only through the
three integer logarithms.  It is therefore constant wherever none of them
changes and moves only where one of them does, and the next definition names
the sets on which they are all constant, so that the next two theorems can say
where the value stands still and by what factor it moves.

Say that $x$ and $y$ lie in the same \emph{logarithmic cell} when
$\lfloor\log_b x\rfloor=\lfloor\log_b y\rfloor$ for each of $b=p,q,r$: the
\lword{Erdos269/ThreePrimeRunningLcm.lean}{163}{SameThreePrimeLogCell}{cell
relation}.

\begin{theorem}[constancy on cells]\label{res:cell}
If $x,y\ge1$ lie in the same logarithmic cell, then
$\Lprefix(x)=\Lprefix(y)$.  The same holds for the kernel at two smooth points
whose values lie in one cell.
\end{theorem}

\begin{proof}
Immediate from Proposition~\ref{res:lcm}, since $\hgt$ depends on $x$ only
through the three integer logarithms.
\end{proof}

\noindent Formalised as the
\lword{Erdos269/ThreePrimeRunningLcm.lean}{179}{smoothPrefixLcm_eq_of_sameLogCell}{cell
constancy of the running value}, the
\lword{Erdos269/ThreePrimeRunningLcm.lean}{169}{threePrimeHeight_eq_of_sameLogCell}{cell
constancy of the height}, and the
\lword{Erdos269/ThreePrimeRunningLcm.lean}{191}{threePrimeKernelQ_eq_of_sameLogCell}{cell
constancy of the kernel}.  The running value therefore changes only when one of
the three logarithms changes, and the next theorem says by exactly how much.

\begin{theorem}[single-coordinate jump ratios]\label{res:jump}
Let $x,y\ge1$.  If $\lfloor\log_p y\rfloor=\lfloor\log_p x\rfloor+1$ while the
other two logarithms agree, then $\Lprefix(y)=p\,\Lprefix(x)$; and similarly
with $q$ or $r$ in place of $p$.
\end{theorem}

\begin{proof}
By Proposition~\ref{res:lcm} both sides are heights, and one factor of the height
gains one in its exponent while the others are unchanged.
\end{proof}

\noindent Formalised as the
\lword{Erdos269/ThreePrimeRunningLcm.lean}{326}{smoothPrefixLcm_firstLogStep}{first},
\lword{Erdos269/ThreePrimeRunningLcm.lean}{339}{smoothPrefixLcm_secondLogStep}{second},
and
\lword{Erdos269/ThreePrimeRunningLcm.lean}{352}{smoothPrefixLcm_thirdLogStep}{third}
coordinate steps, over the corresponding statements for the height alone, which
need no primality: the
\lword{Erdos269/ThreePrimeRunningLcm.lean}{294}{threePrimeHeight_firstLogStep}{first
height step}, the
\lword{Erdos269/ThreePrimeRunningLcm.lean}{305}{threePrimeHeight_secondLogStep}{second},
and the
\lword{Erdos269/ThreePrimeRunningLcm.lean}{316}{threePrimeHeight_thirdLogStep}{third}.

The jump points are therefore the pure prime powers, and they do not collide
across channels.

\begin{theorem}[jump count]\label{res:count}
Let $n\ge0$.  The set of the first $n$ positive powers of $p$, of $q$, and of
$r$ has exactly $3n$ elements, and adjoining the common origin $1$ gives
exactly $3n+1$.
\end{theorem}

\begin{proof}
The count is routine.  Within one channel the powers $b,b^{2},\ldots,b^{n}$ are distinct because
$b\ge2$.  Across two channels a common value would be a positive power of two
distinct primes, which unique factorisation forbids.  Finally $1$ is not a
positive power of any prime.
\end{proof}

\noindent Formalised as the
\lword{Erdos269/ThreePrimeRunningLcm.lean}{249}{threePrimePositiveJumpSet_card}{positive
jump count} and the
\lword{Erdos269/ThreePrimeRunningLcm.lean}{278}{threePrimeJumpSetWithOrigin_card}{jump
count with the origin}, over the
\lword{Erdos269/ThreePrimeRunningLcm.lean}{208}{positivePrimePowers_card}{channel
cardinality}, the
\lword{Erdos269/ThreePrimeRunningLcm.lean}{229}{positivePrimePowers_disjoint}{channel
disjointness}, and the
\lword{Erdos269/ThreePrimeRunningLcm.lean}{219}{one_not_mem_positivePrimePowers}{exclusion
of the origin}; the channels themselves are the
\lword{Erdos269/ThreePrimeRunningLcm.lean}{204}{positivePrimePowers}{positive
power sets}.  The exponent $0$ is omitted from each channel because it is the
shared initial value, which is why the origin is counted once rather than three
times.

Two channels never meet, so at each positive pure power exactly one of the three
logarithms advances, and by Theorem~\ref{res:jump} the running value is
multiplied there by the prime of that channel.  Reading those multipliers in
increasing order of the pure powers gives the \emph{jump word} of $\Lprefix$:
one letter from $\{p,q,r\}$ for each positive pure power, recording which prime
the value is multiplied by at that point.  At $\{2,3,5\}$ the pure powers in
increasing order are $2,3,4,5,8,9,16,25,27,32,\ldots$, so the jump word begins
\[
 2,\;3,2,\;5,2,\;3,2,\;5,3,2,\;\ldots
\]
where the grouping is the one used in Section~\ref{sec:escape}: each group ends
at a power of two.



\section{The height-fibre normal form}\label{sec:fibre}

Grouping a lattice sum by the value of the running least common multiple turns
it into a sum over heights whose coefficients are multiplicities.  We prove
this exactly, on a finite rectangular box.

Write $\mathcal B(h_p,h_q,h_r)$ for the box of exponent triples with $i\le h_p$,
$j\le h_q$, $k\le h_r$, the
\lword{Erdos269/ThreePrimeRunningLcm.lean}{368}{smoothExponentBox}{exponent
box}, and, for a height $H$, write $F(H)$ for the set of points of the box whose
running height $\hgt(p^{i}q^{j}r^{k})$ equals $H$, the
\lword{Erdos269/ThreePrimeRunningLcm.lean}{377}{smoothHeightFiber}{height
fibre}.  Distinct lattice points can carry the same height, and the coefficients
of the normal form are exactly the sizes of these fibres.

\begin{theorem}[finite normal form]\label{res:fibre}
For every box $\mathcal B=\mathcal B(h_p,h_q,h_r)$,
\[
 \sum_{(i,j,k)\in\mathcal B}\Kernel(i,j,k)
 =\sum_{H}\frac{\#F(H)}{H},
\]
the outer sum ranging over the heights actually attained on $\mathcal B$.
\end{theorem}

\begin{proof}
The identity is a regrouping.  Partition $\mathcal B$ into the fibres of the
height map.  On $F(H)$ every summand is
$1/H$ by definition of the kernel, so the fibre contributes $\#F(H)/H$.
\end{proof}

\noindent Formalised as the
\lword{Erdos269/ThreePrimeRunningLcm.lean}{406}{finiteSmoothKernelSum_groupedByHeight}{height-fibre
normal form}, over the
\lword{Erdos269/ThreePrimeRunningLcm.lean}{384}{smoothHeightFiber_kernel_sum}{fibre
sum}, with the height of a lattice point given by the
\lword{Erdos269/ThreePrimeRunningLcm.lean}{373}{smoothPointHeight}{point
height}.

For a worked instance take $(p,q,r)=(2,3,5)$ and the box $\mathcal B(1,1,1)$,
whose eight points carry the smooth values $1,2,3,5,6,10,15,30$ and the heights
\[
\begin{array}{c|cccccccc}
p^{i}q^{j}r^{k}&1&2&3&5&6&10&15&30\\ \hline
\hgt&1&2&6&60&60&360&360&10800
\end{array}
\]
Six heights occur, two of them twice: the points $5$ and $6$ share the height
$60$, and $10$ and $15$ share the height $360$.  Theorem~\ref{res:fibre} here
reads
\[
 1+\tfrac12+\tfrac16+\tfrac1{60}+\tfrac1{60}+\tfrac1{360}+\tfrac1{360}
 +\tfrac1{10800}
 =1+\tfrac12+\tfrac16+\tfrac2{60}+\tfrac2{360}+\tfrac1{10800}
 =\tfrac{18421}{10800}.
\]
The two coefficients $2$ carry the whole content of the regrouping on this box;
where the heights are pairwise distinct the identity is a relabelling and
nothing more.

Together with Theorem~\ref{res:cell} this is the finite core of the ordered
prime-power jump expansion: the value is constant on cells, the cells are
indexed by heights, and the coefficient of a height is the number of lattice
points it collects.  The infinite shell identity is proved in Section~\ref{sec:actual-orbit}.  Theorem~\ref{res:fibre} is an identity between two finite sums,
and it is stated over the full box rather than the smooth prefix, so it is not
a statement about $\Lprefix$ at a cutoff.

The same module records a one-step map $\tau(b,d,s)=b(s-d)$, the
\lword{Erdos269/ThreePrimeRunningLcm.lean}{487}{tailStateStep}{variable-base
tail step}, with the rewriting
\lword{Erdos269/ThreePrimeRunningLcm.lean}{490}{tailStateStep_eq}{that names its
expanded form}.  No orbit of this map is analysed.



\section{A quadratic bound for smooth exponent shells}\label{sec:shell}

A tail estimate needs to know how many lattice points can share a short
multiplicative interval.  Fix a box $\mathcal B(h_p,h_q,h_r)$ and an interval
$[\lo,\hi)$, and write $\mathcal S$ for the set of exponent triples of that box
whose smooth value lies in that interval, the
\lword{Erdos269/ThreePrimeRunningLcm.lean}{500}{smoothExponentShell}{smooth
exponent shell}.  Every bound below assumes the interval short in the
multiplicative sense, meaning that $\hi$ is at most a stated multiple of $\lo$.

Informally, the next lemma says that an interval whose right endpoint is at
most $b$ times its left endpoint contains at most one of the numbers $b^{a}w$.

\begin{lemma}[uniqueness in a short interval]\label{res:short}
Let $b\ge1$ and suppose $\hi\le b\,\lo$.  If $b^{a}w$ and $b^{a'}w$ both lie in
$[\lo,\hi)$, then $a=a'$.
\end{lemma}

\begin{proof}
If $a<a'$ then $\hi\le b\,\lo\le b\cdot b^{a}w=b^{a+1}w\le b^{a'}w<\hi$, which
is impossible; the case $a>a'$ is symmetric.
\end{proof}

\noindent Formalised as the
\lword{Erdos269/ThreePrimeRunningLcm.lean}{509}{exponent_unique_in_short_interval}{short-interval
uniqueness}.  The hypothesis is that the interval has multiplicative width at
most $b$.

\begin{proposition}\label{res:drop}
Let $\mathcal S$ be the shell of the box $\mathcal B(h_p,h_q,h_r)$ in
$[\lo,\hi)$.  If $\hi\le r\,\lo$ then $\#\mathcal S\le(h_p+1)(h_q+1)$; if
$\hi\le p\,\lo$ then $\#\mathcal S\le(h_q+1)(h_r+1)$.
\end{proposition}

\begin{proof}
Suppose $\hi\le r\,\lo$.  If two triples of $\mathcal S$ agree in their first
two coordinates, Lemma~\ref{res:short} applied with $b=r$ and $w=p^{i}q^{j}$
forces their third coordinates to agree as well.  So the projection forgetting
the third coordinate is injective on $\mathcal S$, and its image lies in a
rectangle with $(h_p+1)(h_q+1)$ points.  The case $\hi\le p\,\lo$ is the same
with the first coordinate projected away.
\end{proof}

\noindent Formalised as the
\lword{Erdos269/ThreePrimeRunningLcm.lean}{585}{smoothExponentShell_card_le_dropThird}{third-coordinate
projection} and the
\lword{Erdos269/ThreePrimeRunningLcm.lean}{543}{smoothExponentShell_card_le_dropFirst}{first-coordinate
projection}.

\begin{theorem}[quadratic multiplicity bound]\label{res:shell}
Let $\mathcal S$ be the shell of the box $\mathcal B(h_p,h_q,h_r)$ in
$[\lo,\hi)$.  Suppose $\hi\le r\,\lo$ and $h_p\le h_q\le h_r$ with
$h_p+h_q+h_r=j$.  Then
\[
 9\,\#\mathcal S\le(j+3)^{2}.
\]
\end{theorem}

\begin{proof}
By Proposition~\ref{res:drop}, $\#\mathcal S\le(h_p+1)(h_q+1)$; under the
sorting hypothesis the two surviving coordinates are the two smallest.  It
therefore
suffices to prove that $a\le b\le c$ with $a+b+c=j$ gives
$9(a+1)(b+1)\le(j+3)^{2}$.  From $a\le b\le c$ we get $a+2b\le j$, so it is
enough that $9(a+1)(b+1)\le(a+2b+3)^{2}$; writing $b=a+d$ with $d\ge0$, the
difference of the two sides is $d(3a+4d+3)\ge0$.
\end{proof}

\noindent Formalised as the
\lword{Erdos269/ThreePrimeRunningLcm.lean}{646}{smoothExponentShell_card_quadratic}{quadratic
shell bound}, over the
\lword{Erdos269/ThreePrimeRunningLcm.lean}{629}{sorted_pair_quadratic}{sorted
quadratic estimate}.  The constant $9$ is the square of the number of
generating primes and appears because the bound is the arithmetic--geometric
comparison for a sum of three sorted coordinates; sorting is a hypothesis, not
a normalisation, since the shell itself is not symmetric in the three bases.
The last inequality of the proof is an equality when $h_p=h_q=h_r$, both sides
then being $9(h_p+1)^{2}$, so no constant larger than $9$ survives that step.

The bound is uniform in $\lo$ and $\hi$ subject to the width condition, and it
is stated for the actual filtered shell rather than for a lattice model of it.
It is an input to a tail estimate and is not itself one: no series is bounded
here.  The estimate is elementary and uses no analytic input on the
distribution of smooth numbers, only the projection of
Proposition~\ref{res:drop}.
